\documentclass[../../main.tex]{subfiles}

\begin{document}

\section{开集、闭集、基与子基}\label{section:ch11-开集闭集基子基}

\begin{definition}[拓扑空间]\label{definition:ch11-拓扑空间}
令 $X$ 为非空集. $X$ 的拓扑 $\mathcal{T}$ 是 $X$ 的称为开集的子集族, 具有以下性质:
\begin{enumerate}[(i)]
\item 全空间 $X$ 和空集 $\varnothing$ 是开的.
\item 任何开集的有限族的交是开的.
\item 任何开集族的并是开的.
\end{enumerate}
非空集 $X$ 与 $X$ 上的拓扑一起, 称为拓扑空间. 对于 $X$ 中的点 $x$, 包含 $x$ 的开集称为 $x$ 的邻域.
\end{definition}

我们有时将拓扑空间记为 $(X,\mathcal{T})$. 通常我们对给定的点集仅有一个拓扑感兴趣, 在这样的情形下我们有时用记号 $X$ 既表示点集又表示拓扑空间 $(X,\mathcal{T})$. 当需要更精确时, 我们直接写出拓扑.

\begin{proposition}\label{proposition:ch11-开集邻域}
拓扑空间 $X$ 的子集 $E$ 是开的当且仅当对 $E$ 的每个点 $x$, 存在包含于 $E$ 的 $x$ 的邻域.
\end{proposition}

\begin{proof}
这立即从邻域的定义和拓扑的性质(开集族的并仍然是开的)得到.
\end{proof}

\begin{example}[度量拓扑]\label{example:ch11-度量拓扑}
考虑度量空间 $(X,\rho)$. 定义 $X$ 的子集 $\mathcal{O}$ 为开的, 若对每个点 $x\in\mathcal{O}$ 存在中心在 $x$ 包含于 $\mathcal{O}$ 的开球. 因此开集是开球族的并. 第 9 章的命题 1 断言开集族是 $X$ 的拓扑. 我们称它为由度量 $\rho$ 诱导的度量拓扑. 作为集合上的度量拓扑的特殊情形, 我们有由 $\mathbf{R}^{n}$ 上的欧氏度量诱导的欧氏拓扑.
\end{example}

\begin{example}[离散拓扑]\label{example:ch11-离散拓扑}
令 $X$ 为任何非空集. 定义 $\mathcal{T}$ 为 $X$ 的所有子集族. 则 $\mathcal{T}$ 称为 $X$ 的离散拓扑. 对于离散拓扑, 每个包含点的集合是该点的一个邻域. 离散拓扑由离散度量诱导.
\end{example}

\begin{example}[平凡拓扑]\label{example:ch11-平凡拓扑}
令 $X$ 为任何非空集. 定义 $\mathcal{T}$ 为 $X$ 的由 $\varnothing$ 和 $X$ 组成的子集族. 则 $\mathcal{T}$ 是 $X$ 的拓扑. 对于平凡拓扑, 一个点的仅有的邻域是整个集合 $X$.
\end{example}

\begin{example}[拓扑子空间]\label{example:ch11-拓扑子空间}
给定拓扑空间 $(X,\mathcal{T})$ 和 $X$ 的非空子集 $E$, 定义 $E$ 的拓扑 $\mathcal{S}$ 为所有形如 $E\cap\mathcal{O}$ 的集合, 其中 $\mathcal{O}$ 属于 $\mathcal{T}$. 我们称拓扑空间 $(E,\mathcal{S})$ 为 $(X,\mathcal{T})$ 的子空间.
\end{example}

在初等分析中我们定义 $\mathbf{R}$ 的子集为开的含义时不需要用到"拓扑"这个词. 在第 1 章, 我们证明了拓扑空间 $\mathbf{R}$ 具有性质: 每个开集是开区间的可数不交族的并. 在一个度量空间中, 每个开集是开球族的并.

在一般拓扑空间中, 区分出一族称为基的开集时常是有用的: 它们是拓扑的基石.

\begin{definition}[基与子基]\label{definition:ch11-基}
对于拓扑空间 $(X,\mathcal{T})$ 和 $X$ 中的点 $x$, $x$ 的邻域族 $\mathcal{B}_{x}$ 称为 $x$ 的拓扑基, 若对 $x$ 的任何邻域 $U$, 存在 $\mathcal{B}_{x}$ 中的集合 $B$ 使得 $B\subseteq U$. 开集族 $\mathcal{B}$ 称为拓扑 $\mathcal{T}$ 的基, 若它在每一点包含拓扑的基.

对于拓扑空间 $(X,\mathcal{T})$, $\mathcal{T}$ 的覆盖 $X$ 的子族 $\mathcal{S}$ 称为拓扑 $\mathcal{T}$ 的子基, 若 $\mathcal{S}$ 的有限子族的交是 $\mathcal{T}$ 的基.
\end{definition}

观察到拓扑的一个子族是该拓扑的基当且仅当每个非空开集是 $\mathcal{B}$ 的子族的并. 一旦拓扑的基事先指定, 拓扑就完全定义: 它由 $\varnothing$ 以及属于基的集合的并组成. 拓扑常通过具体化基来定义. 以下命题叙述了 $X$ 的子集族成为拓扑的基必须具有的性质.

\begin{proposition}\label{proposition:ch11-基判别}
对于非空集 $X$, 令 $\mathcal{B}$ 为 $X$ 的子集族. 则 $\mathcal{B}$ 是一个拓扑基当且仅当
\begin{enumerate}[(i)]
\item $\mathcal{B}$ 覆盖 $X$, 即 $X=\bigcup_{B\in\mathcal{B}}B$.
\item 若 $B_{1}$ 和 $B_{2}$ 属于 $\mathcal{B}$ 且 $x\in B_{1}\cap B_{2}$, 则存在 $\mathcal{B}$ 中的集合 $B$ 使得 $x\in B\subseteq B_{1}\cap B_{2}$.
\end{enumerate}
以 $\mathcal{B}$ 为基的唯一的拓扑由 $\varnothing$ 以及 $\mathcal{B}$ 的子族的并组成.
\end{proposition}

\begin{proof}
假设 $\mathcal{B}$ 具有性质(i)和(ii). 定义 $\mathcal{T}$ 为 $\mathcal{B}$ 的子族的并与 $\varnothing$ 组成的族. 我们宣称 $\mathcal{T}$ 是 $X$ 的拓扑. 的确, 我们从(i)推出集合 $X$ 是 $\mathcal{B}$ 中所有集合的并, 因此它属于 $\mathcal{T}$. 此外, $\mathcal{T}$ 的子族的并也是 $\mathcal{B}$ 的子族的并, 因此属于 $\mathcal{T}$. 剩下来要证明: 若 $\mathcal{O}_{1}$ 和 $\mathcal{O}_{2}$ 属于 $\mathcal{T}$, 则交 $\mathcal{O}_{1}\cap\mathcal{O}_{2}$ 属于 $\mathcal{T}$. 事实上, 令 $x$ 属于 $\mathcal{O}_{1}\cap\mathcal{O}_{2}$. 则存在 $\mathcal{B}$ 中的集合 $B_{1}$ 与 $B_{2}$ 使得 $x\in B_{1}\subseteq\mathcal{O}_{1}$, $x\in B_{2}\subseteq\mathcal{O}_{2}$. 用(ii), 选取 $\mathcal{B}$ 中的 $B_{x}$ 满足 $x\in B_{x}\subseteq B_{1}\cap B_{2}$. 则 $\mathcal{O}_{1}\cap\mathcal{O}_{2}=\bigcup_{x\in\mathcal{O}}B_{x}$, 是 $\mathcal{B}$ 的子族的并. 因此 $\mathcal{T}$ 是以 $\mathcal{B}$ 为基的拓扑. 它是唯一的. 我们将逆命题的证明留作练习.
\end{proof}

一组基确定唯一的拓扑. 然而一般情况下, 一个拓扑有许多基. 例如, 开区间族是 $\mathbf{R}$ 上欧氏拓扑的一个基, 而具有有理端点的开有界区间族也是该拓扑的一个基.

\begin{example}[乘积拓扑]\label{example:ch11-乘积拓扑}
令 $(X,\mathcal{T})$ 和 $(Y,\mathcal{S})$ 为拓扑空间. 在笛卡儿积 $X\times Y$ 上, 考虑由乘积 $\mathcal{O}_{1}\times\mathcal{O}_{2}$ 组成的集族 $\mathcal{B}$, 其中 $\mathcal{O}_{1}$ 在 $X$ 中是开的且 $\mathcal{O}_{2}$ 在 $Y$ 中是开的. 我们把检验 $\mathcal{B}$ 是 $X\times Y$ 上的拓扑的基留作练习. 该拓扑称为 $X\times Y$ 上的乘积拓扑.
\end{example}

\begin{example}\label{example:ch11-子基例}
考虑闭有界区间 $[a,b]$ 作为承袭 $\mathbf{R}$ 的拓扑的拓扑空间. 该空间有形如 $[a,c)$ 或 $(c,b]$ 的区间组成的子基, $a<c<b$.

用命题 2 不难看出, $\mathcal{S}$ 的有限子族的交是一个基.
\end{example}

\begin{definition}[闭包点与闭包]\label{definition:ch11-闭包}
对于拓扑空间 $X$ 的子集 $E$, 点 $x\in X$ 称为 $E$ 的闭包点, 若 $x$ 的每个邻域都包含 $E$ 中的点. $E$ 的闭包点全体称为 $E$ 的闭包且记为 $\overline{E}$.
\end{definition}

显然我们总是有 $E\subseteq\overline{E}$. 若 $E$ 包含它的所有闭包点, 即 $E=\overline{E}$, 集合 $E$ 称为闭的.

\begin{proposition}\label{proposition:ch11-闭包最小}
对于拓扑空间 $X$ 的子集 $E$, 它的闭包 $\overline{E}$ 是闭的. 此外, $\overline{E}$ 是包含 $E$ 的最小闭子集, 若 $F$ 是闭集且 $E\subseteq F$, 则 $\overline{E}\subseteq F$.
\end{proposition}

\begin{proof}
集合 $\overline{E}$ 是闭的, 若它包含所有的闭包点. 令 $x$ 为 $\overline{E}$ 的闭包点. 考虑 $x$ 的邻域 $U_{x}$. 存在点 $x'\in E\cap U_{x}$. 由于 $x'$ 是 $E$ 的闭包点而 $U_{x}$ 是 $x'$ 的一个邻域, 存在点 $x''\in E\cap U_{x}$. 因此 $x$ 的每个邻域包含 $E$ 的点, 故 $x\in\overline{E}$. 因此集合 $\overline{E}$ 是闭的. 显然若 $A\subseteq B$, 则 $\overline{A}\subseteq\overline{B}$, 因此若 $F$ 是闭的且包含 $E$, 则 $\overline{E}\subseteq\overline{F}=F$.
\end{proof}

\begin{proposition}\label{proposition:ch11-开集补集}
拓扑空间 $X$ 的子集 $E$ 是开的当且仅当它的补是闭的.
\end{proposition}

\begin{proof}
首先假定 $E$ 在 $X$ 中是开的. 令 $x$ 为 $X\sim E$ 的闭包点. 则 $x$ 不属于 $E$, 这是因为否则将有 $x$ 的一个邻域包含于 $E$, 因此与 $X\sim E$ 不交. 因此 $x$ 属于 $X\sim E$. 现在假定 $X\sim E$ 是闭的. 令 $x$ 属于 $E$. 则必有 $x$ 的邻域包含于 $E$, 否则 $x$ 的每个邻域将包含 $X\sim E$ 中的点, 因此 $x$ 将是 $X\sim E$ 的闭包点. 由于 $X\sim E$ 是闭的, $x$ 属于 $X\sim E$. 这是一个矛盾.
\end{proof}

由于 $X\sim[X\sim E]=E$, 从前一个命题得出拓扑空间 $X$ 的子集是闭的当且仅当它的补在 $X$ 中是开的. 因此, 根据 De Morgan 等式, 拓扑空间的闭子集族具有以下性质.

\begin{proposition}\label{proposition:ch11-闭集性质}
令 $X$ 为拓扑空间. 空集 $\varnothing$ 和全空间 $X$ 是闭的, $X$ 的闭子集的任何有限族的并是闭的, $X$ 的任何闭子集族的交是闭的.
\end{proposition}

\begin{exercise}\label{exercise:ch11-01}
证明非空集 $X$ 的离散拓扑是一个度量拓扑.
\end{exercise}

\begin{exercise}\label{exercise:ch11-02}
证明集合上的离散拓扑有唯一的基.
\end{exercise}

\begin{exercise}\label{exercise:ch11-03}
关于命题 2, 证明若 $\mathcal{B}$ 是拓扑的基, 则性质(i)和(ii)成立.
\end{exercise}

\begin{exercise}\label{exercise:ch11-04}
令 $\mathcal{T}_{1}$ 和 $\mathcal{T}_{2}$ 是非空集 $X$ 的拓扑. 证明 $\mathcal{T}_{1}=\mathcal{T}_{2}$ 当且仅当存在 $\mathcal{T}_{1}$ 的基 $\mathcal{B}_{1}$ 与 $\mathcal{T}_{2}$ 的基 $\mathcal{B}_{2}$, 在 $X$ 中的每个点 $x$ 以下面方式相联系: 对 $x$ 的每个属于 $\mathcal{B}_{1}$ 的邻域 $N_{1}$, 存在 $x$ 的属于 $\mathcal{B}_{2}$ 的邻域 $N_{2}$ 使得 $N_{2}\subseteq N_{1}$, 且对 $x$ 的每个属于 $\mathcal{B}_{2}$ 的邻域 $N_{2}$, 存在 $x$ 的属于 $\mathcal{B}_{1}$ 邻域 $N_{1}$ 使得 $N_{1}\subseteq N_{2}$.
\end{exercise}

\begin{exercise}\label{exercise:ch11-05}
令 $E$ 为拓扑空间 $X$ 的子集.
\begin{enumerate}[(i)]
\item 点 $x\in X$ 称为 $E$ 的内点, 若存在 $x$ 的邻域包含于 $E$; 内点全体称为 $E$ 的内部且记为 $\operatorname{int}E$. 证明 $\operatorname{int}E$ 总是开的且 $E$ 是开的当且仅当 $E=\operatorname{int}E$.
\item 点 $x\in X$ 称为外点, 若存在 $x$ 的邻域包含于 $X\sim E$; 外点全体称为 $E$ 的外部且记为 $\operatorname{ext}E$. 证明 $\operatorname{ext}E$ 总是开的且 $E$ 是开的当且仅当 $\overline{E}\sim E\subseteq\operatorname{ext}E$.
\item 点 $x\in X$ 称为 $E$ 的边界点, 若 $x$ 的每个邻域包含 $E$ 与 $X\sim E$ 中的点; 边界点全体称为 $E$ 的边界且记为 $\operatorname{bd}E$. 证明: (i) $\operatorname{bd}E$ 总是闭的, (ii) $E$ 是开的当且仅当 $E\cap\operatorname{bd}E=\varnothing$, (iii) $E$ 是闭的当且仅当 $\operatorname{bd}E\subseteq E$.
\end{enumerate}
\end{exercise}

\begin{exercise}\label{exercise:ch11-06}
令 $A$ 和 $B$ 为拓扑空间 $X$ 的子集. 证明: 若 $A\subseteq B$ 则 $\overline{A}\subseteq\overline{B}$. 证明 $\overline{A\cup B}=\overline{A}\cup\overline{B}$ 和 $\overline{A\cap B}\subseteq\overline{A}\cap\overline{B}$.
\end{exercise}

\begin{exercise}\label{exercise:ch11-07}
令 $\mathcal{O}$ 为拓扑空间 $X$ 的开子集. 对于 $X$ 的子集 $E$, 证明 $\mathcal{O}$ 与 $E$ 不交当且仅当它与 $\overline{E}$ 不交.
\end{exercise}

\begin{exercise}\label{exercise:ch11-08}
对非空集 $X$ 的子集族 $\mathcal{S}$, 证明存在 $X$ 上的拓扑 $\mathcal{T}$, 它包含族 $\mathcal{S}$ 且具有性质: 任何包含 $\mathcal{S}$ 的其他拓扑也包含 $\mathcal{T}$; 它是包含 $\mathcal{S}$ 的具有最少个数集合的拓扑.
\end{exercise}

\begin{exercise}\label{exercise:ch11-09}
(Sorgenfrey 直线)证明形如 $[a,b)$(其中 $a<b$)的区间族是实数集 $\mathbf{R}$ 的拓扑的基. 具有该拓扑的实数集 $\mathbf{R}$ 称为 Sorgenfrey 直线.
\end{exercise}

\begin{exercise}\label{exercise:ch11-10}
(Moore 平面)考虑上半平面 $\mathbf{R}^{2,+}=\{(x,y)\in\mathbf{R}^{2}\mid y\geqslant 0\}$. 对满足 $y>0$ 的点 $(x,y)$, 取中心为 $(x,y)$ 且包含于上半平面的通常的欧氏开球作为点 $(x,0)$ 的基本开邻域. 取与在 $(x,0)$ 的实直线相切的上半平面的欧氏开球为基本开邻域. 证明该集族是一个基. 具有该拓扑的集合 $\mathbf{R}^{2,+}$ 称为 Moore 平面.
\end{exercise}

\begin{exercise}\label{exercise:ch11-11}
(Kuratowski 14 子集问题)
\begin{enumerate}[(i)]
\item 令 $E$ 为拓扑空间 $X$ 的子集. 证明重复用补与闭包至多能从 $E$ 得到 $14$ 个不同的集合.
\item 在 $\mathbf{R}^{2}$ 中给出从适当的 $E$ 得到 $14$ 个不同的集合.
\end{enumerate}
\end{exercise}

\end{document}